% This is samplepaper.tex, a sample chapter demonstrating the
% LLNCS macro package for Springer Computer Science proceedings;
% Version 2.21 of 2022/01/12
%
\documentclass[runningheads]{llncs}
\usepackage{vmlmacros}

% for random stuff 
\usepackage[T1]{fontenc}
\usepackage{float}

% as recommended 
\pagestyle{plain}
\usepackage{lineno}
\linenumbers 

% T1 fonts will be used to generate the final print and online PDFs,
% so please use T1 fonts in your manuscript whenever possible.
% Other font encondings may result in incorrect characters.
%

% Used for displaying a sample figure. If possible, figure files should
% be included in EPS format.
%
% If you use the hyperref package, please uncomment the following two lines
% to display URLs in blue roman font according to Springer's eBook style:
%\usepackage{color}
%\renewcommand\UrlFont{\color{blue}\rmfamily}
%\urlstyle{rm}
%
\begin{document}
%
\title{Never Going Back (tracking) Again: Equational Pattern Matching with Efficient Compilation}
%
%\titlerunning{Abbreviated paper title}
% If the paper title is too long for the running head, you can set
% an abbreviated paper title here
%

% \author{First Author\inst{1}\orcidID{0000-1111-2222-3333} \and
% Second Author\inst{2,3}\orcidID{1111-2222-3333-4444} \and
% Third Author\inst{3}\orcidID{2222--3333-4444-5555}}

%

% \authorrunning{F. Author et al.}

% First names are abbreviated in the running head.
% If there are more than two authors, 'et al.' is used.
%
% \institute{Princeton University, Princeton NJ 08544, USA \and
% Springer Heidelberg, Tiergartenstr. 17, 69121 Heidelberg, Germany
% \email{lncs@springer.com}\\
% \url{http://www.springer.com/gp/computer-science/lncs} \and
% ABC Institute, Rupert-Karls-University Heidelberg, Heidelberg, Germany\\
% \email{\{abc,lncs\}@uni-heidelberg.de}}

%
\maketitle              % typeset the header of the contribution
%

\begin{abstract}
    Pattern matching appeals to functional programmers for its expressiveness
    and good cost model, but it fails to express certain computations without
    verbosity. Verbosity can be reduced by using extended forms of pattern
    matching, but extensions are not standardized, and making certain extensions
    like or-patterns and pattern guards cooperate in a large-scale programming
    language is a significant implementation challenge that has not seen success
    until recent years. 
    
    Alternatively to pattern matching, computations can be expressed succinctly
    by using equations, as has been demonstrated by the Verse programming
    language. But such computations may have time and space costs that can be
    hard to predict. As a compromise, we ~propose a new language,~\ensuremath{V^{-}},
    which~uses equations in a limited way that makes time and space costs easy
    to predict. In comparative examples, \ensuremath{V^{-}}~ expresses computations as
    succinctly as pattern matching with popular extensions- namely, side
    conditions, or-patterns, and pattern guards- and just like pattern matching,
    it can be compiled to a decision tree. 
    
    We present formal rules and a Standard ML implementation for the \ensuremath{V^{-}}
    match compiler, and we evaluate our implementation with micro-benchmarks and
    a proof that our compilation algorithm preserves the desirable decision tree
    height invariants. 

\end{abstract}

\section{Introduction}

% fancy pm is nice, there are standard ways but not unified, there
%   are other examples like vc and the ultimate, we’re looking at a compromise.
%   contribution is core calculus that is very small yet quite expressive.

\section{Pattern Matching vs. Equations}
\label{pmandextensions}

% --------------------------------------------
\begin{figure}[h]
\centering
\begin{minipage}{0.45\linewidth}
\tiny
\begin{verbatim}

type standard_shape = Square of float 
                    | Triangle of float * float
                    | Trapezoid of float * float * float

let exclaimTall sh =
  match sh with 
    | Square s -> 
        if s > 100.0 
        then "Wow! That's a sizeable square!"
        else "Your shape is small." 
    | Triangle (_, h) -> 
        if h > 100.0 
        then "Goodness! Towering triangle!"
        else "Your shape is small." 
    | Trapezoid (_, _, h) -> 
        if h > 100.0
        then "Zounds! Tremendous trapezoid!"
        else "Your shape is small." 
\end{verbatim}
\caption*{Without extensions}
\end{minipage}
\hfill
\begin{minipage}{0.45\linewidth}
\tiny
\begin{verbatim}
let exclaimTall sh =
  match sh with 
    | Square s when s > 100.0 ->
        "Wow! That's a sizeable square!"
    | Triangle (_, h) when h > 100.0 ->
        "Goodness! Towering triangle!"
    | Trapezoid (_, _, h) when h > 100.0 -> 
        "Zounds! Tremendous trapezoid!"
    | _ -> "Your shape is small." 
\end{verbatim}
\caption*{With side conditions}
\end{minipage}
\caption{Comparing pattern matching with and without side conditions.}
\end{figure}
\begin{figure}[ht]
\centering
\begin{minipage}{0.47\linewidth}
\tiny
\begin{verbatim}
  let tripleLookup (rho : finitemap) (x : int) =
      match lookup rho x with Some w -> 
        (match lookup rho w with Some y -> 
          (match lookup rho y with Some z -> z
          | _ -> handleFailure x)
        | _ -> handleFailure x)
      | _ -> handleFailure x
\end{verbatim}
\caption*{OCaml version with nested pattern matching}
\end{minipage}
\hfill
\begin{minipage}{0.47\linewidth}
\tiny
\begin{verbatim}
  tripleLookup rho x
     | Just w <- lookup rho x
     , Just y <- lookup rho w
     , Just z <- lookup rho y
     = z
  tripleLookup _ x = handleFailure x
\end{verbatim}
\caption*{Haskell version with pattern guards}
\end{minipage}
\caption{Comparing \texttt{tripleLookup} using pattern guards.}
\label{fig:pmtriplelookup}
% \Description{Side-by-side comparison of tripleLookup implementations in OCaml and Haskell.}
\end{figure}
\label{fig:guardtriplelookup}
\begin{figure}[h]
\centering
\begin{minipage}{0.45\linewidth}
\tiny
\begin{verbatim}
let describeColor c =
  match c with
    | Red -> "Warm color"
    | Orange -> "Warm color"
    | Yellow -> "Warm color"
    | Blue -> "Cool color"
    | Green -> "Cool color"
\end{verbatim}
\caption*{Without extensions}
\end{minipage}
\begin{minipage}{0.45\linewidth}
\tiny
\begin{verbatim}
let describeColor c =
  match c with
    | Red | Orange | Yellow -> "Warm color"
    | Blue | Green -> "Cool color"
\end{verbatim}
\caption*{With or-patterns}
\end{minipage}
\caption{Comparing pattern matching with and without or-patterns.}
\end{figure}

% PM is great. But needs to be extended to be more expressive. Consider how 
% pm without extensions is worse: 

% Examples. 

% This Works, but not unified, hard, etc. 

% How much work to extend? How much work to add to new lang? What's missing? 

% Equations are expressive, 

% Examples. But performance is hard to predict. Cite. 


% P+ as just the vessel for the specific examples of extensions. then show v- has
%       all that and more. then focus on dtran. 




\subsection{Introducing \VMinus\ }
\label{vminus}

    \begin{figure}[H]
      \scriptsize
          \begin{center}
            \begin{bnf}
            $P$ :~\textsf{Programs} ::=
            $\bracketed{d}$ : definition
            ;;
            $d$ :~\textsf{Definitions} ::=
            | $\it{val}\; x\;~\it{=}\;~\expr$ : bind name to expression
            ;;
            $\expr$ :~\textsf{Expressions} ::=
            | $\v$ : literal values 
            | $x, y, z$ : names
            | $\bf{if}\;~\left[\; G\;~\bracketed{\dbar\; G}\;~\right]\;~\bf{fi}$ : if-fi 
            | $K~\bracketed{\expr}$ : value constructor application 
            | $\expr[1]\;~\expr[2]$ : function application 
            | $\lambda x.\;~\expr$ : lambda declaration 
            ;;
            $G$ :~\textsf{Guarded Expressions} ::=  
            $[\vexists{\bracketed{x}}]~\bracketed{g}~\boldsymbol{\rightarrow}\expr$ : names, guards, and body
            ;;
            $g$ :~\textsf{Guards} ::=  
            | $\expr$ : intermediate expression 
            | $\expr[1] =~\expr[2]$ : equation 
            | $ g~\bracketed{; g}~\;~\choice~\; g~\bracketed{; g}$ : choice 
            ;;
            $\v$ : Values ::= $K\bracketed{\v}$ : value constructor application 
            | $\lambda x.\;~\expr$ : lambda value
            ;;
            \end{bnf}
            \medskip
          
          * Desugars to~{$x = e_{1};\; x = e_{2}$}, $x$ fresh. 
          \end{center}
          % \Description{The concrete syntax of VMinus.}
              \caption{\VMinus: Concrete syntax}
              \label{fig:vmsyntax}
              \end{figure}      

        \VMinus takes several key concepts from {\VC}---namely, equations and
        choice---with several key modifications. Like~\VC,~\VMinus has
        equations and choice, but unlike in~\VC, choice only guards computations, 
        there are no multiple results, and all decision-making appears in the 
        \iffibf construct, inspired by Dijkstra~\cite{dijkstra}. 

            \begin{figure}[ht] 
      \begin{minipage}[h]{0.54\linewidth}
        \vmlst 
        \begin{lstlisting}[numbers=none, basicstyle=\tiny, xleftmargin=.2em,
          showstringspaces=false,
          frame=single]
val exclaimTall =~\sh.
  if E s. sh = Square s; (> s) 100 -> 
            Wow! That's A Sizeable Square!  
    [] E w h. sh = Triangle w h; (> h) 100 ->
            Goodness! Towering Triangle!
    [] E b1 b2 h. sh = Trapezoid b1 b2 h; 
     (> h) 100 -> Zounds! Tremendous Trapezoid!
    [] ->  Your Shape Is Small
  fi 
  \end{lstlisting}
      % \Description{exclaimTall in VMinus}
      \subcaption{\tt{exclaimTall} in~\VMinus}
          \label{fig:vmexclaimtall} 
      %   \vspace{4ex}
      \end{minipage}%%
      \begin{minipage}[h]{0.5\linewidth}
        \vmlst 
        \begin{lstlisting}[numbers=none, basicstyle=\tiny, xleftmargin=2em,
                      frame=single]
val tripleLookup =~\rho.~\x.
  if E w y z v1 v2 v3. 
    v1 = (lookup rho) x; v1 = Some w; 
    v2 = (lookup rho) w; v2 = Some y; 
    v3 = (lookup rho) y; v3 = Some z 
       -> z 
    [] -> handleFailure x
  fi 
   \end{lstlisting}
        % \Description{tripleLookup in VMinus}
        \subcaption{\tt{tripleLookup} in~\VMinus}
            \label{fig:vmtriplelookup} 
        \vspace{4ex}
      \end{minipage} 
      \begin{minipage}[h]{\linewidth}
        \vmlst 
        \begin{lstlisting}[numbers=none, basicstyle=\tiny, xleftmargin=9em,
          showstringspaces=false,
          frame=single]
val game_of_token =~\t. 
  if E f. (t = BattlePass f  | (t = ChugJug f | t = TomatoTown f)) -> 
        P (Fortnite f)
  [] E f. (t = HuntersMark f | t = SawCleaver f)                   -> 
        P (Bloodborne ((* 2) f))
  [] E f. (t = MoghLordOfBlood f | t = PreatorRykard f)            -> 
        P (EldenRing ((* 3) f))
  [] -> P (Irrelevant 0)
  fi 
\end{lstlisting}
      % \Description{game_of_token in VMinus}
      \subcaption{\tt{game\_of\_token} in~\VMinus}
          \label{fig:vmgot}
      \vspace{4ex}
      \end{minipage}%% 
      \caption{The functions in~\VMinus have a desirably concise
      implementation, as before.}
  \label{fig:vminusfuncs}
    \end{figure}  

\subsection{Formal Semantics of ${\VMinus}$}


    \begin{table}
\begin{tabular}{ll}
\toprule
    \multicolumn2{l}{\emph{\VMinus Metavariables}}~\\
\midrule
    $\expr$& An expression~\\ 
    $\v,\;~\v'$& A value~\\
    $\fail$& An expression failure~\\
    $\result$& $\v~\vert~\fail$ : expressions produce~\it{results}: values or
    failure.~\\
    \Rho& An environment: $name~\rightarrow {\mathcal{V}}_{\bot}$~\\
    $\rho\{ x~\mapsto y~\} $& An environment extended with name $x$ mapping to $y$~\\
    $\g$& A guard~\\
    $\gs$& Zero or more guards, separated by~\it{;}~\\
    $\G$& A guarded expression~\\
    $\Gs$& Zero or more guarded expressions, separated by~\dbar~\\
    $eq$& An equation~\\ 
    % $\tempstuck$& a temporarily-stuck equation~\\
    $\dagger$& when solving guards is rejected~\\
    $s$& $\rhohat~\mid~\dagger$ : guards produce~\it{solutions}: a
    refined environment~\rhohat\; or rejection\\
    $\mathcal{T}$& A context of all temporarily stuck guards (a sequence)~\\ 
    % $G$& A guarded expression~\\
    %~\uppsidown& Inability to compile to a decision tree; a compile time error~\\
\bottomrule
\end{tabular}    
% \Description{Metavariables used in rules for VMinus.}
\caption{\VMinus metavariables and their meanings}
\label{fig:vmmetavars}
\end{table}


    \subsubsection{Expressions}
    
    \newcommand\GNoTree{\vmrung~\rightsquigarrow~\uppsidown} 
    
    An expression in~\VC evaluates to produce possibly-empty sequence of values,
    where an empty sequence of values is identical to the syntactic form~\fail.
    In~\VMinus, an expression never returns multiple values, but it can~\fail.
    Specifically, in~\VMinus, an expression evaluates to produce a single~\it{result}.
    A result is either a single value $\v$ or~\fail. 
    
    \[\result\;~\rm{::=}\;~\v\;~\vert~\;~\fail~\]
    
    \showvjudgement{Eval}{\vmeval}

    \subsubsection{Names and refinement of environments}

    In~\VMinus, like in~\VC, names can be introduced with the existential
    {$\exists$} before they are given a binding. Bindings are given by equations
    in guards~(\ref{guards}). When a name is introduced with {$\exists$}, it is
    bound in the environment~\Rho\ to $\bot$ (pronounced “bottom”). Success of a
    guard only~\it{refines} the environment; that~is, when
    $\grefine[solution=\Rhoprime]$, we~expect $\rho~\subseteq~\Rhoprime$. The
    definition of $\subseteq$ on environments is given below. 
    
    \begin{align*}
    \rho~\subseteq~\Rhoprime~\text{ when }&\dom\rho  \subseteq~\dom~\Rhoprime\\
    \text{ and } &\forall x~\in~\dom~\rho:~\rho(x)~\subseteq~\Rhoprime(x)
    \end{align*}
    
    \smallskip
        
    \subsubsection{Guards}
    \label{guards}

    For example, the~\VMinus expression {\tt{((}\ttbackslash\tt{x. x) K)}}
    succeeds and returns the value~\tt{K}. The~\VMinus expression~\tt{if~~fi},
    the empty~\tt{if-fi}, always {\fail}s. 

    Like~\VC,~\VMinus has a nondeterministic semantics. Guards are solved in~\VMinus
    similarly to how equations are solved in Verse: the program nondeterministically
    picks one out of a context (\context), attempts to solve it, and moves on. 

In my semantics, this process occurs over a~\it{list} of guards~\gs in a guarded
expression \G: the program picks a guard from~\G, attempts to solve it to refine
the environment or fail, and repeats.~\VMinus can only pick a guard out of \G\
that it \it{knows} it can solve. \it{Knowing} a guard can be solved can be determined
in~\VMinus before code is executed. If~\VMinus can't pick a guard and there are
guards left over, the semantics gets stuck before code is executed. 

\showvjudgement{Solve-Guards}{\vmgs}


The environment $\rho$ maps from a name to a value {\v} or $\bot$. $\bot$ means
a name has been introduced with the existential, $\exists$, but is not yet bound
to a value. Given any such $\rho$, a guard~\g\ eithers refine $\rho$
($\Rhoprime$) or is~\bf{rejected}. We use the metavariable $\dagger$ to
represent rejection, and if any guard in a list is rejected, the entire list of
guards is rejected.

    \showvjudgement{Guard-Refine}{\grefine[solution=\rho']}
    \showvjudgement{Guard-reject}\gfail

  For example, in the~\VMinus expression~\tt{if E x. x = K; x = K2 -> x fi}, the
  existential (in concrete syntax,~\tt{E}) introduces~\tt{x} to~\Rho\ bound to
  $\bot$, producing the environment $\bracketed{x~\mapsto~\bot}$. The guard
  \tt{x = K} successfully unifies~\tt{x} with~\tt{K}, producing the environment
  $\bracketed{x~\mapsto~\tt{K}}$. The guard~\tt{x = K2} attempts to unify~\tt{K}
  with~\tt{K2} and is rejected with $\dagger$. 

  \subsection{Notable rules in the \VMinus semantics}

  Worth discussing in the~\VMinus semantics are those in which a name in \Rho\
  is bound to a value or to $\bot$. Notable among these include
  \textsc{Guard-Name-Exp-Bot}, \textsc{Guard-Name-Exp-Eq},
  \textsc{Guard-Name-Exp-Fail}, \textsc{Guard-Names-Bot-Succ},
  and~\textsc{Guard-Names-Bot-Succ-Rev}. This section discusses each of these
  rules in short detail. The full set of rules for \VMinus is in
  Section~\ref{vmsemantics}. 

  In this section, I use the terms~\it{known} and~\it{unknown} to denote a
  name's status in \Rho. If a name $x$ is bound to a value \v in \Rho, then $x$
  is \it{known}. If $x$ is bound to $\bot$ in \Rho, then $x$ is \it{unknown}. I
  use this terminology again when compiling~\VMinus to \D, since it describes
  the same status in both evaluation in compilation. 

\[
\guardnameexpbot
\]

When $\rho(x) = \bot$ ($x$ is \it{unknown}) and $\vmeval$, the program refines
\Rho\ by making a binding of $x$ to \v. Here,~\it{=} is treated like a
let-binding in ML-like languages. 

\[
  \guardnameexpfail
\]
  
When $\rho(x) = \v$ ($x$ is \it{known}) and $\vmeval$, the program proceeds
without refining the environment, since no new information is gained.
Here,~\it{=} is treated like a \tt{==} in C-like languages. 

\[
\guardnameexpfail
\]

When $\rho(x) = \v'$ ($x$ is \it{known}) and $\vmeval$ and $\v \neq \v'$,
unification fails, so the guarded expression fails. Here,~\it{=} is treated like
a \tt{==} in C-like languages. 

\[
\guardnamesbotsucc
\]

When $\rho(x) = \v$ ($x$ is \it{known}) and $rho(y) = \bot$ ($y$ is
\it{unknown}), the program can always bind $y$ to \v. Here,~\it{=} is treated
like a let-binding in ML-like languages. 

\[
\guardnamesbotsuccrev
\]

This rule is simply an application of \textsc{Guard-Names-Bot-Succ} in reverse. 
    
    \vfilbreak
    
    \subsection{Rules (Big-step Operational Semantics) for~\U, shared by~\VMinus and~\D}
    \label{usemantics1}
    \usemantics 
    \subsection{Rules (Big-step Operational Semantics) specific to~\VMinus}
    \label{vmsemantics}
    \vmsemantics

    % An intriguing alternative to pattern matching exists in~\it{equations}, from
    % the Verse Calculus (\VC), a core calculus for the functional logic
    % programming language~\it{Verse}~\citep{antoy2010functional,
    % hanus2013functional, verse}. 
    
    % Powerful. 

    % But hard to predict. 



    

\subsection{\VMinus compiles to a decision tree in the ~\D language}
\label{d}

Like pattern matching, \VMinus~compiles to a decision tree whose height is 
linear with respect to code size. To formalize decision trees, we use 
\D, a language which generalizes the trees found in~\cite{maranget}. 

The spec of decision trees and match compilation algorithm is presented below. 

  \begin{figure}
    \begin{center}
    \dcsyntax
    \end{center}
    % \Description{The concrete syntax of D.}
    \caption{\D: Concrete syntax}
    \label{fig:dsyntax}
    \end{figure}

      In~\D, as in Maranget's trees, the~\it{test} node extracts all names from a
  value constructor at once for use in subtrees. The compiler is responsible for
  introducing the fresh names used in~\it{test}. The compiler alpha-renames all
  necessary terms before it translates an~\iffibf to a decision tree to ensure
  all names are unique. 

  In~\D, like in~\VMinus, expressions can~\it{fail}, meaning some of~\D's
  syntactic forms like~\it{try-let} and~\it{cmp} have an extra branch that is
  executed if the examined expression fails. 


    \subsection{The~\DTran\ algorithm:~\VMinus $\rightarrow$~\D}


    When presented with an~\iffibf,~\DTran\~invokes~\Compile with context~\ctx.
    \Compile~first desugars the~\iffibf by expanding choice to multiple~\iffibf
    branches with a desugaring function~\ITran: 
    
    \itran{if\;~\dots~\dbar\; gs_{1};\;~\choiceg{gs_{2}}{gs_{3}};\; gs_{4}~\rightarrow e\;~\dbar~\dots~\;fi}
    
    ==
    
    ${if\;~\dots~\dbar\; gs_{1};\; gs_{2}~\rightarrow e~\;\dbar\; gs_{3};\; gs_{4}~\rightarrow e~\;\dbar~\dots~\;fi}$

    With the desugared~\iffibf,~\Compile\ then repeatedly chooses a guarded
    expression $G$ and applies one of the compilation rules below. The rules are
    applied in a nondeterministic order. 

    \Compile terminates when it inserts a final~\it{match} node (with rule
    \textsc{Match}).  A~\it{match} node is inserted for a right-hand side expression~\expr when the list of
    guards preceding~\expr is empty or a list of assignments from names to
    unbound names. Termination of~\DTran\~is guaranteed because each recursive
    call passes a list of guarded expressions in which the number of guards is
    strictly smaller, so eventually the algorithm reaches a state in which the
    first unmatched branch is all trivially-satisfied guards.


    If~\Compile\ cannot choose a $g$ of one of the valid forms, it halts with an
    error. This can happen when no $g$ is currently solvable in the context, as
    determined by the same algorithm that~\VMinus uses to pick a guard to solve,
    or when the program would be forced to unify incompatible values, such as
    any value with a closure. 

    \rab{I have these rules with english descriptions, though that increases length}.

    
\showvjudgement{Compile}{\compilebig}

\subsection{Big-step rules for \Compile}

      \begingroup
      \scriptsize
      \compiler
      \endgroup

      Application of the rules is nondeterministic, which is desirable in a
      match compiler: the structure of the final result~\expr\ depends on the
      \textit{heuristic input} to the compiler, which we discuss in
      section~\ref{reductionstrats}.


\section{Implementation}

\section{Evaluation and Analysis}

\subsection{Microbenchmarks}


\subsection{Reduction Strategies}
\label{reductionstrats}

      In my implementation, I apply the~\textsc{Test} rule before other rules,
      having found in my experiments that this heuristic leads to the smallest trees. A~%
      risk of inserting a \it{test} node before a \it{let-unless} node is that
      if there are many shared names across branches, a \it{test} node
      will introduce those names in \it{let-unless} nodes in each
      subtree.
      It is possible to insert \it{let-unless} nodes before \it{test}
      nodes, in which case those names will be bound
      only once.  The risk of the \it{let}-first strategy
      is that if there are many names used in only one subtree of a \it{test},
      these names will be introduced to \it{all} branches.  The extra
      bindings may harm
      performance by bloating an environment in an interpreter or thrashing the
      icache if the tree further is compiled to machine code. 

      In the implementation,~\DTran\~first introduces all the names under the
      all existential $\exists$'s to a context which determines if a name is
      \it{known} or~\it{unknown}, applying the~\textsc{Exists} rule is applied
      all at the start. At the start of compilation, each name introduced by
      $\exists$ is~\it{unknown} in a context~\ctx. Since all names in the
      program are unique at this stage, there are no clashes. 


\subsection{Proof Sketch: \Compile~Preserves the Decision Tree Invariant}

\section{Related and Future Work}
    


    \section{Conclusion}

    % We~have introduced \VMinus, a~language that makes decisions using equations.
    % By~example, we have shown that programs written in~\VMinus can have the same
    % desirable properties as equivalent programs written using pattern matching.
    % And~to~show that equations can be compiled to efficient decision trees, we
    % have introduced~\D and~\DTran. In doing so, we~have demonstrated that
    % programming with equations is a promising alternative to pattern matching.
    
    

\begin{credits}
\subsubsection{\ackname} A bold run-in heading in small font size at the end of the paper is
used for general acknowledgments, for example: This study was funded
by X (grant number Y).

\subsubsection{\discintname}
It is now necessary to declare any competing interests or to specifically
state that the authors have no competing interests. Please place the
statement with a bold run-in heading in small font size beneath the
(optional) acknowledgments\footnote{If EquinOCS, our proceedings submission
system, is used, then the disclaimer can be provided directly in the system.},
for example: The authors have no competing interests to declare that are
relevant to the content of this article. Or: Author A has received research
grants from Company W. Author B has received a speaker honorarium from
Company X and owns stock in Company Y. Author C is a member of committee Z.
\end{credits}
%
% ---- Bibliography ----
%
% BibTeX users should specify bibliography style 'splncs04'.
% References will then be sorted and formatted in the correct style.
%
% \bibliographystyle{splncs04}
% \bibliography{mybibliography}
%
\begin{thebibliography}{8}
\bibitem{ref_article1}
Author, F.: Article title. Journal \textbf{2}(5), 99--110 (2016)

\bibitem{ref_lncs1}
Author, F., Author, S.: Title of a proceedings paper. In: Editor,
F., Editor, S. (eds.) CONFERENCE 2016, LNCS, vol. 9999, pp. 1--13.
Springer, Heidelberg (2016). \doi{10.10007/1234567890}

\bibitem{ref_book1}
Author, F., Author, S., Author, T.: Book title. 2nd edn. Publisher,
Location (1999)

\bibitem{ref_proc1}
Author, A.-B.: Contribution title. In: 9th International Proceedings
on Proceedings, pp. 1--2. Publisher, Location (2010)

\bibitem{ref_url1}
LNCS Homepage, \url{http://www.springer.com/lncs}, last accessed 2023/10/25
\end{thebibliography}
\end{document}